\input texsis
\paper

\overfullrule=0pt
\def\frac#1#2{#1\over #2}
%
\def\bmatrix#1{\left[ \matrix{#1} \right]}
\def\be{{\beta}}
\def\toone{{t+1}}

%\def\frac#1#2{#1\over #2}
%\magnification=1200
\overfullrule=0pt

\line{PhD Macroeconomics \hfil NYU}
\line{October 12, 2023 \hfil Thomas Sargent}
%
% \noindent Macroeconomics Q1 \hfill NYUAD \\
% Fall, 2019 \hfill Thomas Sargent
% \vspace{.3cm}

\centerline{\bf Homework 3}
%\centerline{\bf Homework 3}


\medskip

\noindent{\bf Instructions:}  Please feel free to work in groups of 3 or 4. But please hand in  your own work
before the end of the TA session on Friday October 20.

%\noindent {\bf 1. 15 points  \quad}

\medskip
\noindent{\bf 1.  \ } 
  A consumer has preferences over consumption streams
that are ordered by
the utility functional
$$  E_0 \sum_{t=0}^\infty \beta^t u(c_t)  \leqno(1) $$
where $E_t$ is the mathematical expectation conditioned
on the consumer's time $t$ information,  $c_t$ is time $t$ consumption,
$u(c)$ is  utility function, and
$\beta \in (0,1)$ is a discount factor.  The consumer maximizes
(1) by choosing a consumption, borrowing plan
 $\{c_t, b_{t+1}\}_{t=0}^\infty$ subject to the sequence of budget constraints
$$ c_t + b_t = R^{-1} b_{t+1}  + y_t  , \quad t \geq 0$$
where $\{y_t\}$ is an exogenous
 stationary endowment process, $R$ is a constant gross
risk-free interest rate, $b_t$ is one-period risk-free  bond   maturing at
$t$, and $b_0 $ is a given initial debt that is due at time $0$.  Assume that for $t \geq 1 $, $b_t$ is chosen at time $t-1$.  Assume
that $R^{-1} = \beta$.  Assume that the endowment
process is governed by
$$ \eqalign{ y_{t+1} & = \alpha_c
+ \sum_{j=1}^2 \rho_j y_{t-j+1}  + \sum_{j=1}^2
 \delta_j  z_{t+1-j}  + \psi_1 w_{1,t+1} \cr
  z_{t+1} & = \sum_{j=1}^2 \gamma_j y_{t-j+1} + \sum_{j=1}^2 \phi_j z_{t-j+1}
+ \psi_2 w_{2,t+1}
   \cr} \leqno(1)   $$
where $\{w_{t+1}\}$ is an i.i.d.  sequence of  $2 \times 1$ vectors of random shock whose time $t+1$ component is  distributed as ${\cal N}(0,I)$
%martingale difference
%sequence,  adapted to \hfil \break $J_t = \left[\matrix{w_t & \ldots w_1 &
%c_{0} & c_{-1} & z_0 & z_{-1}\cr}\right]$, with contemporaneous
%covariance matrix $E w_{t+1}w_{t+1}' | J_t = I$, 
 and the
coefficients  $\rho_j, \delta_j, \gamma_j, \phi_j$ are such
that
the matrix
$$ A = \left[\matrix{ \rho_1 & \rho_2 & \delta_1 &\delta_2 & \alpha_c \cr
                      1 & 0 & 0 & 0 & 0 \cr
                     \gamma_1 & \gamma_2 & \phi_1 &\phi_2 & 0 \cr
                     0 & 0 & 1 & 0 & 0 \cr
                     0 & 0 & 0 & 0 & 1 \cr }\right]$$
has eigenvalues less than or equal to unity in modulus.

\medskip
\noindent {\bf a.}. Assume that the utility function $u(c) = \ln c$ %satisfies $u' > 0, u''< 0$ 
for all $c \geq 0$ and  the ``borrowing limit'' $b_t \leq \bar b$ for some $\bar b > 0$. Formulate the
consumer's problem as a dynamic program.
Write pseudo-code to solve it.
\medskip
\noindent{\bf b.}
Now  assume the
the quadratic  utility function
$$u(c_t) =  -.5 (c_t - \gamma)^2 \leqno(2) $$
where $\gamma >0$ is a bliss level of consumption.
Assume that $\bar b =+\infty$ so that there is no borrowing limit.
%
%GGHH 
%We assume that the endowment
%process has the state-space representation
%$$ \eqalign{ z_{t+1} & = \check A z_t + \check C w_{t+1}\cr
%               y_t & = \check G  z_t \cr}$$
%where $w_{t+1}$ is an i.i.d.~process with mean zero and
%identity contemporaneous covariance matrix, $\check A$ is a stable matrix,
%its eigenvalues being strictly below unity in modulus, and
%$\check G$ is a vector. The {\it state\/} vector confronting the household at $t$  is
%$\left[\matrix{b_t & z_t\cr}\right]'$, where $b_t$ is its one-period debt falling
% due at the beginning of period $t$
%and $z_t$ contains all variables useful for
%forecasting its future endowment.
Formulate the consumer's problem as a dynamic program and guess a solution.
Justify your guess.

\medskip
\noindent {\bf c.}. Continue  to assume utility function (2), but now replace the part {\bf a}   borrowing limit  with 
the following condition on the
consumption, borrowing plan:
$$  % \lim_{T \rightarrow +\infty}
 E_0 \sum_{t=0}^\infty \beta^t b_t^2 < +\infty.  $$
%This condition suffices to rule out Ponzi schemes.
Formulate the consumer's problem as a dynamic program and write pseudo-code to solve it.


\medskip


\medskip
\noindent {\it \bf 3.}  \quad 
A worker who  lives during periods $t = 0, 1, 2, \ldots, + \infty$  has an opportunity to work in period $t$ at wage $w_t \geq 0 $ and disutility of labor $\xi_t \geq 0$.  Each period a previously unemployed worker has the opportunity to draw a  new wage from a c.d.f. $F$ that satisfies $F(0) = 0, F(B) = 1$ for $B >0$.  In period $t$, a worker   employed last period at wage $w_{t-1}$ has the opportunity  not to draw again but instead to retain  $w_{t-1} $  and to  set $w_t = w_{t-1}$.  But if the worker quits, she cannot draw a new offer this period and in the future
she cannot recall that offer or any other  past offer.  A worker draws one offer from $F$ each period she starts a period after being unemployed the period before.  An employed  worker receives time $t$ utility $w_t - \xi_t$, where each period $\xi_t$ is a nonnegative random variable drawn afresh from c.d.f. $G$ that satisfies $G(0) = 0, G(D)= 1$.  Successive draws of $\xi_t$ from $G$ are independently and identically distributed.
Successive new draws of $w$ from $F$ are independently and identically distributed. A worker receives unemployment benefits
$b \geq 0 $ each period she is unemployed. Where $\beta \in (0,1)$ and time $t$ utility is
$$u_t = \cases{w_t - \xi_t&
if working  \cr
\Big. $b$  & if not working, \cr} $$
  a worker wants to maximize  
the mathematical expectation of 
$
\sum_{t=0}^\infty \beta^t u_t 
$
by sequentially deciding whether  to retain last period's job being employed  at wage $w_{t-1}$ or  quit and become unemployed and receive $b$ for a period and then next period draw a new wage offer $w$ and  disutility of working $\xi$. 

\medskip
\noindent {\bf a.}
Please tell how the following functional equation    generates  a statistical model of quits and exits from unemployment:
$$
\eqalign{ v(w_t, \xi_t) & = \max_{{\rm accept}, {\rm reject}} \left\{ w_t - \xi_t + \beta p(w_t) , b + \beta Q \right\} \cr
    p(w_t) & = \int v(w_t, \xi') d G (\xi') \cr
    Q & = \int v(w, \xi) d F(w) d G(\xi) . } \leqno(1)
$$
Please carefully describe the objects $v(w, \xi), p(w),$ and $Q$.

\medskip
\noindent{\bf b.} Please describe a previously employed worker's optimal policy for quitting.
\medskip

\noindent{\bf c.} Please describe a previously unemployed worker's optimal policy for accepting a new job.
\medskip
\noindent{\bf d.} \quad {\it Extra Credit} Assume that $G$ and $F$ are both uniform distributions.  Write Python, Julia, 
or Matlab code to solve functional equation (1) when $b = B/4$,  $D= B/2$, and $\beta = .95$. Describe the worker's
optimal decisions quantitatively.  
\bye

\bye

